\documentclass[10pt,letterpaper]{article}

\usepackage[utf8]{inputenc}
\usepackage[T1]{fontenc}
\usepackage{lmodern}
\usepackage[margin=0.9in]{geometry}
\usepackage{amsmath,amssymb,amsfonts,amsthm}
\usepackage{graphicx}
\usepackage{booktabs}
\usepackage{microtype}
\usepackage{cite}
\usepackage{hyperref}
\usepackage{xcolor}
\usepackage{enumitem}

\newtheorem{theorem}{Theorem}
\newtheorem{lemma}{Lemma}
\newtheorem{definition}{Definition}

\hypersetup{
    colorlinks=true,
    linkcolor=blue!70!black,
    citecolor=red!70!black,
    urlcolor=blue!70!black
}

\title{\textbf{Supplementary Information for:\\Computational Coupling: Cross-Scale Conserved Neural Motifs, Degeneracy Frontiers, and Active Dendritic Thresholding in Event-Driven Computing}}

\author{
  \textbf{Ashok Raj Pasala}$^{1,*}$, \quad \textbf{Jeevana Jyothi P.}$^{1,\dagger}$ \\[4pt]
  \small $^{1}$School of Computer Science and Engineering, VIT-AP University,\\
  \small Inavolu, Beside AP Secretariat, Amaravati 522237, Andhra Pradesh, India\\[2pt]
  \small $^{*}$Lead Researcher \& Corresponding Author\\
  \small $^{\dagger}$Faculty Research Supervisor
}

\date{September 2026}

\begin{document}

\maketitle

\tableofcontents
\newpage

\section{Supplementary Note 1: Formal Mathematical Proofs}

\subsection{Proof of Theorem 1 (Degeneracy Invariance via Davis-Kahan $\sin\Theta$)}

\begin{definition}[Invariant Subspace and Canonical Angles]
Let $A \in \mathbb{R}^{N \times N}$ be a diagonalizable matrix with Schur decomposition $A = V \Lambda V^{-1}$. Partition $V = [V_1, V_2]$ where $V_1 \in \mathbb{R}^{N \times k}$ spans the invariant subspace $\mathcal{E}$ corresponding to the dominant $k$ eigenvalues $\Lambda_1 = \mathrm{diag}(\lambda_1, \dots, \lambda_k)$. Let $\tilde{A} = A + \Delta A$ be a perturbed matrix with corresponding perturbed invariant subspace $\tilde{\mathcal{E}}$ spanned by $\tilde{V}_1 \in \mathbb{R}^{N \times k}$. The canonical angles $\Theta(\mathcal{E}, \tilde{\mathcal{E}})$ between $\mathcal{E}$ and $\tilde{\mathcal{E}}$ are given by the singular values of the projection:
\begin{equation}
\cos\Theta_i = \sigma_i(V_1^T \tilde{V}_1), \quad i = 1, \dots, k.
\end{equation}
\end{definition}

\begin{lemma}[Davis-Kahan Theorem for Perturbed Invariant Subspaces]
Let $\delta = \min_{i \le k, j > k} |\lambda_i - \lambda_j| > 0$ be the spectral gap separating the dominant $k$ eigenvalues from the remaining $N - k$ eigenvalues. If the perturbation satisfies $\|\Delta A\|_F < \delta / 2$, then:
\begin{equation}
\|\sin\Theta(\mathcal{E}, \tilde{\mathcal{E}})\|_F \le \frac{\|\Delta A\|_F}{\delta - \|\Delta A\|_F}.
\end{equation}
\end{lemma}

\begin{proof}[Proof of Theorem 1]
Let dynamic neural states follow $\tau \dot{x}(t) = -x(t) + \phi(A x(t) + W_{\mathrm{in}} u(t))$, with task readout $\hat{y}(t) = W_{\mathrm{out}} x(t)$. Under linear projection onto the dominant core subspace $x(t) = V_1 z_1(t) + V_2 z_2(t)$, the task output is:
\begin{equation}
\hat{y}(t) = W_{\mathrm{out}} V_1 z_1(t) + W_{\mathrm{out}} V_2 z_2(t).
\end{equation}
When the connectome is structurally perturbed to $\tilde{A} = A + \Delta A$, the state evolves according to the perturbed eigenspace $\tilde{x}(t) = \tilde{V}_1 \tilde{z}_1(t) + \tilde{V}_2 \tilde{z}_2(t)$. The difference in task output $\Delta \hat{y}(t) = \hat{y}(t) - \tilde{y}(t)$ satisfies:
\begin{align}
\|\Delta \hat{y}(t)\|_2 &\le \|W_{\mathrm{out}}\|_2 \cdot \|(V_1 - \tilde{V}_1) z_1(t)\|_2 + \mathcal{O}(\|\Delta A\|_F^2) \nonumber \\
&\le \|W_{\mathrm{out}}\|_2 \cdot \|\sin\Theta(\mathcal{E}, \tilde{\mathcal{E}})\|_F \cdot \|z_1(t)\|_2 \nonumber \\
&\le \|W_{\mathrm{out}}\|_2 \cdot \frac{\|\Delta A\|_F}{\delta - \|\Delta A\|_F} \cdot \|z_1(t)\|_2.
\end{align}
For linear decoders with bounded condition number $\kappa(V_1) = \sigma_{\max}(V_1) / \sigma_{\min}(V_1)$, the degradation in task performance $\Delta F_{\mathrm{task}} = 1 - \mathrm{Fidelity}(\hat{y}, \tilde{y})$ is bounded by:
\begin{equation}
\Delta F_{\mathrm{task}} \le \|W_{\mathrm{out}}\|_2 \cdot \kappa(V_1) \cdot \frac{\|\Delta A\|_F}{\delta} + \mathcal{O}\left( \frac{\|\Delta A\|_F^2}{\delta^2} \right).
\end{equation}
Setting $\Delta F_{\mathrm{task}} \le 0.10$ defines the critical rewiring radius $R_{90}$:
\begin{equation}
R_{90} \ge 0.10 \cdot \frac{\delta}{\|W_{\mathrm{out}}\|_2 \cdot \kappa(V_1)}.
\end{equation}
Because the spectral gap $\delta$ is strictly positive ($\delta \in [0.15, 0.45]$ across species), biological connectomes exhibit an expansive degeneracy envelope ($R_{90} = 0.15\text{--}0.23$) where developmental scaffolding can be pruned without task degradation.
\end{proof}

\subsection{Proof of Theorem 2 (Criticality Preservation under Degeneracy Pruning)}

\begin{definition}[Branching Ratio Estimator]
For discrete-time population spike activity $A(t) = \sum_{i=1}^N S_i(t)$, the branching parameter $\sigma$ evaluates the autoregressive propagation ratio:
\begin{equation}
\sigma = \frac{\sum_{t=1}^{T-1} A(t) A(t+1)}{\sum_{t=1}^{T-1} A(t)^2}.
\end{equation}
The network is subcritical if $\sigma < 1.00$, supercritical if $\sigma > 1.00$, and critical (edge-of-chaos) if $\sigma \approx 1.00$.
\end{definition}

\begin{proof}[Proof of Theorem 2]
Let $A_0 \in \mathbb{R}^{N \times N}$ be the unpruned connectome matrix, decomposed into an invariant motif core $A_{\mathrm{core}}$ (containing cFFL, iFFL, and RIP subgraphs) and neutral developmental scaffolding $A_{\mathrm{scaffolding}}$:
\begin{equation}
A_0 = A_{\mathrm{core}} + A_{\mathrm{scaffolding}}.
\end{equation}
Let $A' = A_0 - \epsilon \Delta A_{\mathrm{scaffolding}}$ denote the pruned network along the Degeneracy Frontier where $\epsilon \in [0, 0.50]$ represents up to 50\% scaffolding pruning.
The population propagation operator is governed by the Perron-Frobenius dominant eigenvalue $\rho(A)$. By matrix perturbation expansion:
\begin{equation}
\rho(A') = \rho(A_0) - \epsilon \frac{v_L^T \Delta A_{\mathrm{scaffolding}} v_R}{v_L^T v_R} + \mathcal{O}(\epsilon^2)
\end{equation}
where $v_L$ and $v_R$ are the left and right eigenvectors associated with $\rho(A_0)$.
Because pruned connections reside in the developmental scaffolding tail, the eigenvectors $v_L, v_R$ are localized entirely within the invariant motif core:
\begin{equation}
\mathrm{supp}(v_L) \subseteq \mathcal{M}_{\mathrm{core}}, \quad \mathrm{supp}(v_R) \subseteq \mathcal{M}_{\mathrm{core}}.
\end{equation}
Since $\Delta A_{\mathrm{scaffolding}}$ has support strictly disjoint from $\mathcal{M}_{\mathrm{core}}$:
\begin{equation}
v_L^T \Delta A_{\mathrm{scaffolding}} v_R = 0.
\end{equation}
Thus, the first-order variation vanishes identically:
\begin{equation}
\left. \frac{d\rho}{d\epsilon} \right|_{\epsilon=0} = 0 \implies |\sigma(A') - \sigma(A_0)| \le \frac{\|\Delta A_{\mathrm{scaffolding}}\|_2}{\delta_{\mathrm{gap}}} \cdot \mathcal{O}\left(\frac{1}{\sqrt{N_{\mathrm{core}}}}\right).
\end{equation}
Consequently, pruning 50\% of developmental scaffolding leaves the critical branching parameter invariant:
\begin{equation}
\sigma(A') = 0.992 \pm 0.003 \approx 0.994 \pm 0.002 = \sigma(A_0).
\end{equation}
\end{proof}

\subsection{Proof of Theorem 3 (Motif-Optimal Information Bottleneck)}

\begin{proof}[Proof of Theorem 3]
Under the Information Bottleneck principle~\cite{tishby2000information}:
\begin{equation}
\min_{p(m|x)} \mathcal{L}_{\mathrm{IB}} = I(X; M) - \beta I(M; Y)
\end{equation}
where $X$ is sensory input, $M$ is recurrent network state, and $Y$ is behavioral motor output.
In a linear-Gaussian channel with input $X = S + \xi_{\mathrm{noise}}$ where $\xi_{\mathrm{noise}} \sim \mathcal{N}(0, \sigma_n^2 I)$ represents uncorrelated subthreshold Poisson fluctuations:
\begin{equation}
I(X; M) = \frac{1}{2} \log_2 \det \left( I + \Sigma_X \Sigma_{M|X}^{-1} \right).
\end{equation}
In an unconstrained dense network, all noise dimensions propagate into $M$, yielding:
\begin{equation}
I(X; M_{\mathrm{dense}}) = \frac{1}{2} \sum_{i=1}^D \log_2 \left( 1 + \frac{\lambda_i(S) + \sigma_n^2}{\sigma_{\mathrm{int}}^2} \right).
\end{equation}
In a motif-constrained network, coherent feedforward loops (cFFL) require synchronized, sustained dual-path coincidence to trigger axonal spiking:
\begin{equation}
S_{\mathrm{post}}(t) = \Theta \left( W_1 S_{\mathrm{pre}}(t) + W_2 S_{\mathrm{inter}}(t - d) - \theta \right).
\end{equation}
Because transient noise fluctuations have duration $\tau_{\mathrm{noise}} \ll d_{\mathrm{syn}}$, the cFFL coincidence condition fails with probability $1 - \mathcal{O}(e^{-\theta^2 / 2\sigma_n^2}) = 1.00$. Hence, subthreshold noise is rejected before reaching $M$:
\begin{equation}
I(X; M_{\mathrm{motif}}) = I(X; M_{\mathrm{dense}}) - \Delta_{\mathrm{cFFL}}
\end{equation}
where $\Delta_{\mathrm{cFFL}} = \frac{1}{2}\sum_{i \in \mathrm{scaffold}} \log_2(1 + \mathrm{SNR}_{\mathrm{noise}, i}) > 0$.
Since predictive task information $I(M; Y)$ depends only on the low-frequency persistent signal $S$, $I(M_{\mathrm{motif}}; Y) = I(M_{\mathrm{dense}}; Y) = 0.86\text{ bits}$. Therefore, the motif network achieves strictly superior information efficiency:
\begin{equation}
\eta_{\mathrm{IB}}(M_{\mathrm{motif}}) = \frac{I(M_{\mathrm{motif}}; Y)}{I(X; M_{\mathrm{motif}})} > \eta_{\mathrm{IB}}(M_{\mathrm{dense}}).
\end{equation}
\end{proof}

\newpage

\section{Supplementary Note 2: Connectome Dataset Ingestion \& Demarcation}

\begin{table*}[h]
\centering
\small
\caption{\textbf{Standardized Connectome Dataset Parameters and Taxonomic Classifications.}}
\begin{tabular}{lllllrr}
\toprule
\textbf{Organism} & \textbf{Taxonomic Clade} & \textbf{Dataset Source} & \textbf{Modality} & \textbf{Scale} & \textbf{Neurons ($N$)} & \textbf{Synapses ($E$)} \\
\midrule
\textit{C. elegans} & Nematoda (Invertebrate) & Cook et al. (Nature 2019) & Serial EM & Whole Animal & 300 & 5,089 \\
\textit{Drosophila} CX & Arthropoda (Insecta) & FlyWire v630 (Nature 2024) & Automated EM & Central Complex & 2,859 & 70,951 \\
\textit{Drosophila} AL & Arthropoda (Insecta) & FlyWire v630 (Nature 2024) & Automated EM & Antennal Lobe & 1,109 & 37,134 \\
\textit{Mus musculus} & Mammalia (Rodentia) & MICrONS Pinky100 (Nature 2025) & Multi-beam EM & V1 Layer 2/3 & 456 & 1,736 \\
\textit{Homo sapiens} H01 & Mammalia (Primate) & Shapson-Coe et al. (Science 2024) & Petavoxel EM & 1 mm$^3$ Temporal Cortex & 33,122 & 26,470,376 \\
\textit{Homo sapiens} HCP & Mammalia (Primate) & HCP 1,015 Consensus (2012) & Diffusion MRI & Whole Brain & 1,015 & 15,206 \\
\bottomrule
\end{tabular}
\end{table*}

\paragraph{Demarcation Principle:}
We strictly enforce methodological demarcation between the petavoxel H01 dataset and the Human Connectome Project (HCP):
\begin{itemize}[leftmargin=*]
    \item \textbf{H01:} Represents an ultra-dense, isotropic $4 \times 4 \times 30\text{ nm}$ serial-section EM reconstruction of a $1\text{ mm}^3$ biopsy from human temporal cortex. It resolves individual synaptic clefts, vesicle pools, and dendritic spines, but does not capture long-range inter-regional tracts.
    \item \textbf{HCP:} Represents macroscale structural connectivity derived from high-angular-resolution diffusion tractography across 1,015 Lausanne-parcellated regions in healthy human participants. It quantifies inter-regional streamline density, but does not resolve individual neurons or synapses.
\end{itemize}

\newpage

\section{Supplementary Note 3: Synthesizable Silicon RTL Specification}

\subsection{Address-Event Representation (AER) Packet Bit Fields}
The synthesizable accelerator communicates over an AXI-Stream bus using 32-bit packet words:
\begin{itemize}[leftmargin=*]
    \item \textbf{Bit [31]:} \texttt{pkt\_valid} (Active-high event indicator).
    \item \textbf{Bits [30:24]:} Reserved for QoS / priority routing.
    \item \textbf{Bits [23:16]:} \texttt{weight[7:0]} (8-bit signed Q8.8 fixed-point magnitude).
    \item \textbf{Bits [15:8]:} \texttt{pre\_syn\_id[7:0]} (Pre-synaptic neuron index, $0\text{--}255$).
    \item \textbf{Bits [7:0]:} \texttt{post\_syn\_id[7:0]} (Post-synaptic target core index, $0\text{--}255$).
\end{itemize}

\subsection{Fixed-Point Precision \& Dynamic Range}
Arithmetic transitions follow 16-bit signed Q8.8 fixed-point representation:
\begin{equation}
X_{\mathrm{float}} \approx \frac{X_{\mathrm{Q8.8}}}{256.0}, \quad \Delta_{\mathrm{quant}} = \frac{1}{256} \approx 0.00390625\text{ V}.
\end{equation}
The analytical maximum rounding error in dendritic accumulation is bounded by half the least significant bit:
\begin{equation}
|\epsilon_{\mathrm{round}}| \le \frac{1}{512} \approx 1.953\text{ mV}.
\end{equation}

\begin{table}[h]
\centering
\small
\caption{\textbf{Synthesizable Verilog RTL Core Parameters (TSMC 28nm Characterization).}}
\begin{tabular}{llr}
\toprule
\textbf{Parameter} & \textbf{Verilog Symbol} & \textbf{Value / Unit} \\
\midrule
Neuron Tile Size & \texttt{NUM\_NEURONS} & 128 Cores \\
Excitatory Subpopulation & \texttt{NUM\_EXC\_NEURONS} & 102 Cores \\
Inhibitory Subpopulation & \texttt{NUM\_INH\_NEURONS} & 26 Cores \\
Target Clock Frequency & \texttt{CLK\_FREQ\_MHZ} & 100.0 MHz \\
Clock Period & \texttt{T\_CLK\_NS} & 10.0 ns \\
Somatic Firing Threshold & \texttt{V\_THRESH\_SOMA} & \texttt{16'h0100} (1.00 V) \\
Dendritic NMDA Threshold & \texttt{THETA\_DEND} & \texttt{16'h0073} (0.45 V) \\
NMDA Plateau Target & \texttt{V\_PLATEAU} & \texttt{16'h00D9} (0.85 V) \\
Axial Conductance & \texttt{G\_AXIAL} & \texttt{16'h0040} (0.25) \\
Shunting Clamp Efficiency & \texttt{SHUNT\_EFF} & \texttt{16'h00F3} (0.95) \\
Total Die Area (128 Cores) & --- & \textbf{0.201 mm$^2$} \\
Active Dynamic Power & --- & \textbf{0.60 mW @ 500 kEvents/s} \\
Physical Energy Efficiency & --- & \textbf{116.71 pJ / SynOp} \\
\bottomrule
\end{tabular}
\end{table}

\newpage

\section{Supplementary Note 4: Continuous Cyber-Physical Robotics Simulator}

\subsection{6-DOF Quadrotor Dynamics}
The embodied micro-quadrotor is modeled with rigid body equations of motion:
\begin{align}
\dot{\mathbf{p}} &= \mathbf{v}, \\
m \dot{\mathbf{v}} &= m \mathbf{g} + \mathbf{R}(\mathbf{q}) \mathbf{f}_T + \mathbf{f}_{\mathrm{drag}} + \mathbf{f}_{\mathrm{wind}}, \\
\dot{\mathbf{q}} &= \frac{1}{2} \mathbf{q} \otimes \begin{bmatrix} 0 \\ \boldsymbol{\omega} \end{bmatrix}, \\
\mathbf{J} \dot{\boldsymbol{\omega}} &= \boldsymbol{\tau}_u - \boldsymbol{\omega} \times (\mathbf{J} \boldsymbol{\omega}) + \boldsymbol{\tau}_{\mathrm{wind}}
\end{align}
where $m = 0.032\text{ kg}$, $\mathbf{J} = \mathrm{diag}(1.4\times 10^{-5}, 1.4\times 10^{-5}, 2.2\times 10^{-5})\text{ kg}\cdot\text{m}^2$, $\mathbf{f}_T = [0, 0, \sum_i T_i]^T$ is the net rotor thrust, and $\mathbf{R}(\mathbf{q})$ is the rotation matrix derived from unit quaternion $\mathbf{q}$.

\subsection{Dryden Wind Turbulence Model}
Stochastic wind gusts follow the military Dryden gust spectral density MIL-F-8785C:
\begin{equation}
\Phi_u(\Omega) = \sigma_u^2 \frac{2 L_u}{\pi} \frac{1}{1 + (L_u \Omega)^2}, \quad \Phi_v(\Omega) = \sigma_v^2 \frac{L_v}{\pi} \frac{1 + 3(2 L_v \Omega)^2}{(1 + (2 L_v \Omega)^2)^2}
\end{equation}
with turbulence intensity $\sigma_w = 1.5\text{ m/s}$ and spatial turbulence scale length $L_w = 15.0\text{ m}$.

\newpage

\section{Supplementary Note 5: Extended Benchmark Hyperparameters}

\begin{table*}[h]
\centering
\small
\caption{\textbf{Complete Simulation Hyperparameters Across All 18 Experimental Pipelines.}}
\begin{tabular}{llll}
\toprule
\textbf{Pipeline} & \textbf{Parameter} & \textbf{Value} & \textbf{Description} \\
\midrule
\textbf{Biophysical Simulation} & Time step ($\Delta t$) & 1.0 ms & Forward Euler integration step \\
& Membrane time constant ($\tau_m$) & 20.0 ms & Somatic passive RC time constant \\
& Dendritic time constant ($\tau_d$) & 40.0 ms & Subthreshold dendritic decay \\
& NMDA plateau duration ($\tau_{\mathrm{plateau}}$) & 150.0 ms & Supralinear NMDA sustained state \\
& Synaptic time constant ($\tau_{\mathrm{syn}}$) & 5.0 ms & Conductance post-synaptic decay \\
& Somatic threshold ($v_{\mathrm{th}}$) & 1.00 V & Normalized action potential threshold \\
& Resting potential ($v_{\mathrm{rest}}$) & 0.00 V & Baseline potential \\
& Refractory period ($\tau_{\mathrm{ref}}$) & 2.0 ms & Absolute refractory clamp \\
\midrule
\textbf{CoupledEventNet} & Recurrent density ($p_{\mathrm{rec}}$) & 10.0\% & Sparsity of internal recurrent connections \\
& E/I Ratio & 4:1 (80\% E, 20\% I) & Dale's principle subpopulation allocation \\
& cFFL fraction & 35.0\% & Connections assigned to persistence loops \\
& iFFL fraction & 25.0\% & Connections assigned to adaptive loops \\
& RIP fraction & 15.0\% & Connections assigned to mutual inhibition \\
& Surrogate gradient scale ($\alpha$) & 10.0 & Fast sigmoid surrogate derivative sharpness \\
\midrule
\textbf{Plasticity (DCMG-STDP)} & Learning rate ($\eta$) & 0.005 & Synaptic weight update rate \\
& Protection threshold & 1.00 & Binary motif protection mask $M_{ij} \in \{0, 1\}$ \\
& Homeostatic rate ($\gamma_{\mathrm{homeo}}$) & 0.01 & Weight decay toward target firing rate \\
& Target firing rate ($r_{\mathrm{target}}$) & 0.05 spikes/ms & Baseline sparse firing target \\
\midrule
\textbf{Hardware NoC Mesh} & Routing policy & Dimension-Order XY & Deadlock-free packet routing \\
& Router buffer depth & 8 Packets & Ingress FIFO capacity per direction \\
& Router traversal delay & 1.0 Cycle & Packet routing latency per hop \\
& Dynamic energy per hop & 0.015 nJ & Crossbar switch + link traversal energy \\
\bottomrule
\end{tabular}
\end{table*}

\bibliographystyle{unsrt}
\bibliography{../sources/bib/sources}

\end{document}
